\documentclass[UTF8,11pt]{ctexart}
\usepackage[a4paper,margin=25mm]{geometry}
\usepackage{amsmath,amssymb,amsthm,mathtools}
\usepackage[all]{xy}
\usepackage{xurl}
\usepackage{hyperref}
\usepackage{enumitem}
\hypersetup{hidelinks,breaklinks=true}
\setlength{\emergencystretch}{3em}
\newtheorem*{theorem}{Theorem}
\newtheorem*{lemma}{Lemma}
\newtheorem*{proposition}{Proposition}
\newtheorem*{corollary}{Corollary}
\newtheorem*{definition}{Definition}
\newcommand{\Spec}{\operatorname{Spec}}
\newcommand{\Hom}{\operatorname{Hom}}
\newcommand{\End}{\operatorname{End}}
\newcommand{\Gal}{\operatorname{Gal}}
\newcommand{\Cl}{\operatorname{Cl}}
\newcommand{\vol}{\operatorname{vol}}
\newcommand{\covol}{\operatorname{covol}}
\newcommand{\tr}{\operatorname{tr}}
\newcommand{\im}{\operatorname{im}}
\newcommand{\id}{\operatorname{id}}
\newcommand{\coker}{\operatorname{coker}}
\newcommand{\stacksref}[2]{\href{https://stacks.math.columbia.edu/tag/#1}{#2 [#1]}}
\begin{document}
\section*{001\quad Serre仿射性判据}
\subsection*{来源与整理范围}
\textbf{作者：} The Stacks Project Authors。\textbf{作品：} \emph{The Stacks Project}, Chapter 30, \emph{Cohomology of Schemes}, Section 30.3, Lemma 30.3.1, Tag 01XF。

\textbf{原文：} \url{https://stacks.math.columbia.edu/tag/01XF}。
原生 TeX：\url{https://github.com/stacks/stacks-project/blob/master/coherent.tex}，原文件第 282--382 行（本次读取版本）；Git blob SHA：\texttt{d0e7f49ad29ce3092e40c766d630b9ad99d8e1ba}。

\textbf{取得日期：} 2026-09-28。网页原命题及完整证明已取得；另对照了该章 PDF 第 4 页。所用原生 TeX 题证摘录保存在 \texttt{sources/stacks/01XF\_original\_excerpt.tex}。当前下载通道未能保存整章原始文件，不把摘录称为整章原件。

\textbf{原文文献指向：} \href{https://stacks.math.columbia.edu/bibliography/Serre-criterion}{Serre-criterion}；\href{https://stacks.math.columbia.edu/bibliography/EGA}{EGA}, II, Theorem 5.2.1 (d'), IV (1.7.17)。


\subsection*{作者原命题与人类证明}
\begin{lemma}[30.3.1; Serre's criterion for affineness]
Let $X$ be a scheme.
Assume that
\begin{enumerate}
\item $X$ is quasi-compact,
\item for every quasi-coherent sheaf of ideals
$\mathcal{I} \subset \mathcal{O}_X$ we have $H^1(X, \mathcal{I}) = 0$.
\end{enumerate}
Then $X$ is affine.
\end{lemma}

\begin{proof}
Let $x \in X$ be a closed point. Let $U \subset X$ be an affine open
neighbourhood of $x$. Write $U = \Spec(A)$ and let
$\mathfrak m \subset A$ be the maximal ideal corresponding to $x$.
Set $Z = X \setminus U$ and $Z' = Z \cup \{x\}$.
By Schemes, Lemma \stacksref{01J3}{26.12.4} there
are quasi-coherent sheaves of ideals
$\mathcal{I}$, resp.\ $\mathcal{I}'$ cutting out
the reduced closed subschemes $Z$, resp.\ $Z'$.
Consider the short exact sequence
$$
0 \to \mathcal{I}' \to \mathcal{I} \to \mathcal{I}/\mathcal{I}' \to 0.
$$
Since $x$ is a closed point of $X$ and $x \not \in Z$ we see that
$\mathcal{I}/\mathcal{I}'$ is supported at $x$. In fact, the restriction
of $\mathcal{I}/\mathcal{I'}$ to $U$ corresponds to the $A$-module
$A/\mathfrak m$. Hence we see that $\Gamma(X, \mathcal{I}/\mathcal{I'})
= A/\mathfrak m$. Since by assumption $H^1(X, \mathcal{I}') = 0$
we see there exists a global section $f \in \Gamma(X, \mathcal{I})$
which maps to the element $1 \in A/\mathfrak m$ as a section of
$\mathcal{I}/\mathcal{I'}$. Clearly we have
$x \in X_f \subset U$. This implies that $X_f = D(f_A)$ where
$f_A$ is the image of $f$ in $A = \Gamma(U, \mathcal{O}_X)$.
In particular $X_f$ is affine.

\medskip\noindent
Consider the union $W = \bigcup X_f$ over all $f \in \Gamma(X, \mathcal{O}_X)$
such that $X_f$ is affine. Obviously $W$ is open in $X$.
By the arguments above every closed point of
$X$ is contained in $W$. The closed subset $X \setminus W$ of $X$
is also quasi-compact
(see Topology, Lemma \stacksref{005C}{5.12.3}).
Hence it has a closed point if it is nonempty (see
Topology, Lemma \stacksref{005E}{5.12.8}).
This would contradict the fact that all closed points are in
$W$. Hence we conclude $X = W$.

\medskip\noindent
Choose finitely many $f_1, \ldots, f_n \in \Gamma(X, \mathcal{O}_X)$
such that $X = X_{f_1} \cup \ldots \cup X_{f_n}$ and such that each
$X_{f_i}$ is affine. This is possible as we've seen above.
By Properties, Lemma \stacksref{01QF}{28.28.3}
to finish the proof it suffices
to show that $f_1, \ldots, f_n$ generate the unit ideal in
$\Gamma(X, \mathcal{O}_X)$. Consider the short exact sequence
$$
\xymatrix{
0 \ar[r] &
\mathcal{F} \ar[r] &
\mathcal{O}_X^{\oplus n} \ar[rr]^{f_1, \ldots, f_n} & &
\mathcal{O}_X \ar[r] &
0
}
$$
The arrow defined by $f_1, \ldots, f_n$ is surjective since the
opens $X_{f_i}$ cover $X$. We let $\mathcal{F}$ be the kernel
of this surjective map.
Observe that $\mathcal{F}$ has a filtration
$$
0 = \mathcal{F}_0 \subset \mathcal{F}_1 \subset
\ldots \subset \mathcal{F}_n = \mathcal{F}
$$
so that each subquotient $\mathcal{F}_i/\mathcal{F}_{i - 1}$ is
isomorphic to a quasi-coherent sheaf of ideals.
Namely we can take $\mathcal{F}_i$ to be the intersection of
$\mathcal{F}$ with the first $i$ direct summands of
$\mathcal{O}_X^{\oplus n}$.
The assumption
of the lemma implies that $H^1(X, \mathcal{F}_i/\mathcal{F}_{i - 1}) = 0$
for all $i$. This implies that
$H^1(X, \mathcal{F}_2) = 0$ because it is sandwiched between
$H^1(X, \mathcal{F}_1)$ and $H^1(X, \mathcal{F}_2/\mathcal{F}_1)$.
Continuing like this we deduce that $H^1(X, \mathcal{F}) = 0$.
Therefore we conclude that the map
$$
\xymatrix{
\bigoplus\nolimits_{i = 1, \ldots, n} \Gamma(X, \mathcal{O}_X)
\ar[rr]^{f_1, \ldots, f_n} & &
\Gamma(X, \mathcal{O}_X)
}
$$
is surjective as desired.
\end{proof}


\subsection*{整理说明（非作者正文）}
保留作者命题、英文证明、论证次序和两个交换图；仅把外部 \texttt{\string\ref} 改为可定位的章节号与 Tag 链接，并将原注记中的名称移入标题。不添加逆命题，不把相邻的有限型版本另计一题。

记号：$X_f$ 是整体函数 $f$ 非零的开子概形；$D(f_A)$ 是仿射谱中的主开集。标准先修为概形、拟凝聚层、层上同调长正合列及正文指向的基本仿射判据；本题自身承担的是由 $H^1$ 消失构造整体函数、覆盖与单位理想的论证。选题难度依据这一构造与上同调转换，属于高级代数几何证明，不以网页称其为 lemma 或课程名称作难度判断。

\subsection*{可修改的简短标注（非作者正文）}
\textbf{$c$：} 概形的仿射性。

\textbf{$a$：} 拟凝聚理想层；层上同调；短正合列；整体截面提升；层的滤过；主开覆盖与单位理想。

\texttt{cross\_theory\_status = yes}。
把仿射性问题转化为理想层短正合列上的 $H^1$ 消失与整体截面提升，再把所得截面变回仿射主开覆盖和单位理想。位置：原证明第 1 段及最后一段（核层 $\mathcal F$ 的滤过）。以上为非穷尽、可修改初稿。
\subsection*{许可与修改记录}
原数学正文作者：The Stacks Project Authors（集体作者）。原文按 GNU Free Documentation License 1.2 或后续版本发布；本摘录的整理部分随原文采用相同许可。许可全文随包保存在 \texttt{sources/stacks/GFDL-1.2.txt}。无不变章节、封面文字或封底文字。本文题名与来源作品题名不同，不暗示原作者认可本次选题或标注。

\textbf{History.} 2026-09-28：ChatGPT 为本对话材料包整理摘录，加入题号、来源定位、独立排版与简短标注；数学正文的具体排版变动见整理说明。

\end{document}
